\section{Introduction}
\label{sec:intro}

Neural representations of quantum wavefunctions are now used both as
variational ans\"atze for many-body problems \citep{carleo2017,hermann2020,pfau2020}
and as mesh-free solvers for the stationary Schr\"odinger equation
\citep{raissi2019,karniadakis2021}. In the second setting a network $\psi_\theta$ is trained so
that $\Hh\psi_\theta\approx E\psi_\theta$ in a discrete $L^2$ sense, and the
paper reports an energy error and a wavefunction overlap against a reference
solution.

This evaluation scheme has a structural weakness. The reference is either an
analytic solution---available only for textbook potentials---or the output of
a classical solver on the same grid. In the second case the comparison
measures agreement between two approximations, not accuracy, and it silently
absorbs the discretization error of the grid into the "method error". More
importantly, in a genuinely new problem there is no reference at all, and the
only quantities available are those the approximation itself provides. A
residual is such a quantity, but making it small is itself delicate---the
gradient pathologies of residual training are well documented
\citep{wang2021}---and even a small residual is not an error: a residual of
$10^{-2}$ may correspond to an energy error of $10^{-6}$ or to a wavefunction
that is entirely the wrong eigenstate. Both directions occur in the runs
reported below. A network whose energy was accurate to seven digits
($|\eta-\lambda_h|=2.6\times10^{-7}$ on the infinite-well grid) carried a
residual of $0.53$, so that the certificate could prove only an angle bound
of $0.56$~rad and an enclosure six orders of magnitude wider than the true
error (Section~\ref{sec:leak-results}). Conversely, a network trained with a
scale-invariant loss returned a field whose residual norm was
$2.9\times10^{-4}$ while its $L^2$ distance from the eigenfunction of unit norm
was $0.99$: the residual had been made small by shrinking the field
(\ref{app:failures}).

The numerical analysis literature solved the corresponding problem for finite
element and spectral methods long ago: eigenvalue enclosures follow from
variational bounds (Weinstein, Kato--Temple, Lehmann--Goerisch) together with
a verified count of the spectrum in an interval \citep{babuska1991,parlett1998},
and eigenvector bounds follow from the $\sin\theta$ theorems of
Davis and Kahan \citep{daviskahan1970}. Verified-computing packages routinely
return such enclosures with directed rounding. What has not been done, to our
knowledge, is to connect this machinery to neural eigenpairs, where the
"approximation" is a mesh-free function that must first be sampled on a grid,
where the hypothesis of the classical theorems (an interval containing exactly
one eigenvalue) is not available a priori, and where the training procedure
itself can leak the reference solution into the object being evaluated.

This paper closes that gap with a protocol and a reference implementation. The
ingredients are classical; the contribution is that they are assembled into
something that can be run on a neural eigenpair and that verifies its own
hypotheses. Specifically:

\begin{enumerate}
\item \textbf{A computable certificate} (Section~\ref{sec:cert}) that takes an
  approximate eigenfunction $\psi$ on a grid and returns
  (i) an interval $[a,b]$ that is \emph{verified} to contain exactly one
  eigenvalue of the discrete Hamiltonian $\Hh$,
  (ii) a two-sided enclosure of that eigenvalue from the Kato--Temple bound,
  and (iii) an upper bound on the angle between $\psi$ and the exact discrete
  eigenfunction from the Davis--Kahan theorem. The verification of (i) uses
  Sturm sequences in one dimension and the Sylvester inertia of a no-pivot
  sparse LU factorization in two; it never uses a reference eigenvalue.
\item \textbf{Two numerical safeguards} that the verification needs in
  floating point: an explicit check that the factorization performed no
  pivoting, and a \emph{resolution floor} below which an interval is declared
  unverifiable. The second is what makes the single-state protocol refuse a
  degenerate state instead of certifying it with a spuriously narrow interval
  (Section~\ref{sec:floor}).
\item \textbf{A cluster certificate} (Section~\ref{sec:subspace}) for the case
  the single-state certificate must refuse. At a degenerate level no
  eigenfunction is defined by the problem, so we certify the invariant
  subspace instead: the multiplicity is measured by the same verified count,
  each cluster eigenvalue is bracketed to within the residual of its Ritz
  pair, and a principal-angle bound (Theorem~\ref{thm:angle}) certifies the
  subspace spanned by a set of network outputs without separating their Ritz
  values. The construction also yields the failure modes that matter in
  practice---a subspace that missed a state is refused because its dimension
  does not match the count, and a set of independent training runs that
  collapsed onto the same representative is either refused (when the only
  second direction they contain is the difference of two nearly identical
  outputs) or accepted with a bound orders of magnitude weaker than a genuine
  pair's.
\item \textbf{A protocol} (Section~\ref{sec:protocol}) that uses the
  certificate in four ways: to remove leakage from the training residual by
  evaluating the energy at the Rayleigh quotient of the current iterate, to
  select the truncation order of a Dirichlet spectral filter a posteriori
  without any reference, and to carry the certified angle bound back to the
  unfiltered network output through an elementary transport inequality, and
  to decide which sets of independent outputs actually span a degenerate
  level (Section~\ref{sec:deg-results}).
\item \textbf{A directed-rounding count} (Section~\ref{sec:certified}): the
  one-dimensional Sturm count, and the two-dimensional one for separable
  potentials, are recomputed in interval arithmetic with outward rounding, so
  that a verified count is a theorem about the stored matrix rather than a
  guarded floating-point estimate. An unresolved pivot or an unresolved sum
  is \emph{refused}, with its position, never rounded to an integer; bisection
  on the certified count returns certified brackets for discrete eigenvalues
  in one and two dimensions.
\item \textbf{A variational bound on the continuous operator}
  (Section~\ref{sec:variational}): any conforming trial function is a witness
  that $\lambda_1(\mathcal H)$ does not exceed its Rayleigh quotient, and the
  P1/Q1 interpolant of a sampled eigenfunction---including a network's---is
  conforming, so the bound on $\mathcal H$ costs one quadrature and no
  reference solution, no grid operator and no interpolation constant. The
  matching lower bound needs a comparison principle; we obtain it for the
  infinite well and the box and show it to fail for the finite well, whose
  certified discrete energy can lie above the continuous one.
\item \textbf{An empirical audit} (Section~\ref{sec:results}) on six
  potentials in one and two dimensions, three seeds each, including a
  degeneracy experiment, a controlled splitting that locates the boundary
  between state and subspace certification, a leakage comparison, and a
  separation of network error from discretization error.
\end{enumerate}

We state clearly what is \emph{not} claimed. We do not propose new spectral
inequalities: Weinstein's, Temple's and Davis--Kahan's bounds are classical,
and the cluster statements of Section~\ref{sec:subspace} are rearrangements of
the same two arguments applied to an invariant subspace rather than to a
vector; what is new is the construction that makes them \emph{decidable}
for a degenerate level, where the single-state hypothesis cannot hold. The two-sided
enclosures of Sections~\ref{sec:cert}--\ref{sec:subspace} refer to the
finite-difference operator $\Hh$, and the bridge to the continuum in
Section~\ref{sec:certified} is one-sided: an upper bound for every potential,
a lower bound only where a comparison principle exists. Interval arithmetic
does not reach the general two-dimensional count either---the separable case
is exact, while the interval elimination refuses on most shifts, and we report
that measurement rather than route around it. Finally, we do not
claim that the certificate is tight. It is not, for non-smooth potentials; we
report the ratio between certified width and true error rather than hiding it.

The intended use of the protocol is therefore not to replace an eigensolver.
On every problem in this paper a classical solver produces the eigenpair in a
fraction of the time, and the certificate applied to that eigenpair is
$10^{7}$--$10^{10}$ times narrower than the one applied to a trained network,
except on the two infinite wells, where the filtered network is itself the
exact sine mode and the two certificates coincide at the round-off floor.
The intended use is to answer, for a network that has already been trained,
the question "how wrong can this be?", and to make that answer an auditable
number in a table.
